% ===================================================================
%  TAULUKKO N:o 13 — Hotelli. — Ravintola. — Kahvila.
%  Traduction 2026 d'apres texto/io, controlee sur texto/fr et
%  texto/sv. Consigne : voir texto/fi/10-taulukko-01.tex.
%
%  Deux notes, marquees toutes deux « (*) », et deux appels : celui du
%  bloc 01-2 (la chambre) et celui du bloc 09-1 (le menu).
%
%  UNE INVERSION EVITEE D'AVANCE, au bloc 05-1 : « la konduktisto (30)
%  dil omnibuso (29) ». Relative, comme toujours.
% ===================================================================

%%K t13-noto-1 noto
\VUnotes{143.2mm}{%
\textit{\VUgras{(*) Katso huoneen yksityiskohdista taulukko n:o 6.}}%
}

%%K t13-apar-1 apar
\VUsaut{3.0mm}
\VUtitre{15.0pt}{60}{TAULUKKO N:o 13}
\VUsaut{1.6mm}
\VUpk{10.2pt}{40}{Hotelli. --- Ravintola.}
\VUsaut{2.0mm}
\VUpk{10.2pt}{40}{Kahvila.}
\VUsaut{3.4mm}

%%K t13-01-1 p
1. --- Tultuani eilen illalla Royaniin asetuin
« \textit{Valtameren Hotelliin} », jota eräs ystävä oli suositellut minulle.

%%K t13-01-2 p
Minulle annettiin kaunis \VUgras{huone} \textsuperscript{(2)} toisessa
\VUgras{kerroksessa} \textsuperscript{(1)}, jossa nukuin hyvin hyvin (*).

%%K t13-02-1 p
2. --- Tänä aamuna, kun menin ulos huoneestani, johon \VUgras{palvelustyttö}
\textsuperscript{(3)} oli kantanut aamiaisen, menin \VUgras{tupakkahuoneeseen}
\textsuperscript{(5)}, jota käytetään myös \VUgras{lukuhuoneena}
\textsuperscript{(4)}, lukeakseni \VUgras{sanomalehtiä} \textsuperscript{(6)},
samalla kun poltin \VUgras{savukkeen} \textsuperscript{(8)}. Lukemisen jälkeen
laskeuduin \VUgras{hissillä} \textsuperscript{(9)} ja menin kävelemään kaupungin
läpi.

%%K t13-03-1 p
3. --- Koska olin unohtanut kysyä hotellin \VUgras{virkailijalta}
\textsuperscript{(12)}, mihin aikaan tarjoillaan ateriat \VUgras{yhteisessä
pöydässä} \textsuperscript{(11)}, palasin aikaisin ja käytin sen mennäkseni
kylpemään hotellin \VUgras{kylpyhuoneeseen} \textsuperscript{(10)}. Sitten menin
taas \VUgras{kassalle} \textsuperscript{(15)} soittaakseni eräälle ystävälleni.
Mutta \VUgras{puhelin} \textsuperscript{(13)} oli varattu; kun
\VUgras{kassanhoitajatar} \textsuperscript{(14)}, joka kirjoitti \VUgras{laskua}
\textsuperscript{(17)} \VUgras{matkustajarekisteriin} \textsuperscript{(16)},
näki hämmennykseni, hän pyysi \VUgras{portieeria} \textsuperscript{(18)}
lähettämään \VUgras{hotellipojan} \textsuperscript{(19)} hoitamaan asiani
\VUgras{polkupyörällä} \textsuperscript{(20)}.

%%K t13-04-1 p
4. --- Kiitin häntä ja menin ulos antaakseni juomarahan \VUgras{kantajalle}
\textsuperscript{(24)} (tavarankantajalle), joka asemalta
\VUgras{käsikärryllään} \textsuperscript{(25)} oli tuonut \VUgras{hattulaatikkoni}
\textsuperscript{(26)}, \VUgras{matkalaukkuni} \textsuperscript{(27)} ja
\VUgras{matkakassini} \textsuperscript{(28)}. \VUgras{Kirjeenkantaja}
\textsuperscript{(21)}, joka oli juuri tullut, antoi minulle \VUgras{kirjeen}
\textsuperscript{(22)}.

%%K t13-05-1 p
5. --- Viivyin vielä jonkin aikaa oven kynnyksellä \VUgras{postilaatikon}
\textsuperscript{(23)} luona katsellakseni liikennettä kadulla.
\VUgras{Konduktööri} \textsuperscript{(30)}, joka on \VUgras{omnibussissa}
\textsuperscript{(29)}, lastasi vaunuunsa \VUgras{matka-arkun}
\textsuperscript{(31)}, joka kuului \VUgras{matkustajalle}
\textsuperscript{(32)}, jolle \VUgras{hotellinisäntä} \textsuperscript{(33)}
sanoi jäähyväiset. Nuori \VUgras{sähkösanomanlähetti} \textsuperscript{(35)}
kantoi kiirehtimättä \VUgras{sähkettä} \textsuperscript{(34)} eräälle hotellin
vieraalle; \VUgras{turisti} \textsuperscript{(36)}, \VUgras{kamera}
\textsuperscript{(38)} olallaan, tutki \VUgras{matkaopastaan}
\textsuperscript{(37)}.

%%K t13-tit-1 sub
\VUsaut{4.0mm}
\VUpk{10.2pt}{40}{Aamiaisen aika.}
\VUsaut{3.0mm}

%%K t13-06-1 p
6. --- Ristikkoaidan edessä torkkui huoleton \VUgras{kaupunginlähetti}
\textsuperscript{(39)} \VUgras{kantokoukkunsa} \textsuperscript{(40)} ja
\VUgras{kiillotuslaatikkonsa} \textsuperscript{(41)} välissä, samalla kun nuori
\VUgras{pesijätär} \textsuperscript{(42)}, joka kantoi \VUgras{koria}
\textsuperscript{(43)} täynnä liinavaatteita, nauroi \VUgras{hovimestarin}
\textsuperscript{(44)} juhlalliselle ilmeelle, joka seisoi ovella.

%%K t13-07-1 p
7. --- Kun näin \VUgras{keittäjän} \textsuperscript{(45)}, joka tuli ulos
\VUgras{keittiöstään} \textsuperscript{(47)} \VUgras{kellarikerroksessa}
\textsuperscript{(48)} lähettääkseen \VUgras{tiskipojan}
\textsuperscript{(46)} hoitamaan asiaa, muistin, että aamiaisen hetki on
lähellä, ja menin ravintolaan kulkiessani pihan yli, jolla
\VUgras{autonkuljettaja} \textsuperscript{(50)} pani sisään \VUgras{autonsa}
\textsuperscript{(49)}, samalla kun \VUgras{mekaanikko} \textsuperscript{(52)}
korjasi \VUgras{pyöräilijän} \textsuperscript{(53)} ohjeiden mukaan
\VUgras{petrolikolmipyörää} \textsuperscript{(51)}, jolle oli juuri sattunut
onnettomuus.

%%K t13-08-1 p
8. --- Kysyin nuorelta \VUgras{hotellipojalta} \textsuperscript{(54)}, joka
kantoi \VUgras{olutlaseja} \textsuperscript{(55)}, onko yhteinen pöytä kuistin
alla, ja hänen myöntävän vastauksensa jälkeen otin paikan yhdessä pöydistä ja
annoin antaa itselleni erinomaisen aterian.

%%K t13-noto-2 noto
\VUnotes{143.2mm}{%
\textit{\VUgras{(*) Sanavarastoon liitetyn ruokalistan avulla opettaja voi
helposti muunnella allaolevaa ruokalistaa.}}%
}

%%K t13-tit-2 sub
\VUsaut{4.0mm}
\VUpk{10.2pt}{40}{Erinomainen aamiainen.}
\VUsaut{3.0mm}

%%K t13-09-1 p
9. --- Tarjoilija toi minulle aamiaisen ruokalistan (*) ja viinilistan. Valitsin
ja tilasin \VUgras{alkuruoaksi katkarapuja} \VUgras{voin} kera, ja
väliruoaksi \VUgras{komagea Caenin tapaan}. En syönyt \VUgras{kalaa}, vaan
pyysin annoksen \VUgras{lampaanpaistia}, ja \VUgras{vihannesruoaksi pinaattia}
kerman kera. Sitten, koska yksikään \VUgras{väliruoista} ei miellyttänyt minua,
annoin tuoda erilaisia jälkiruokia, \VUgras{juustoa}, \VUgras{hedelmiä} ja
\VUgras{korppuja}. Lisäsin sitä paitsi erinomaista Saint-Émilion-\VUgras{viiniä},
ja hyvin tyytyväisenä ateriaani jätin pöydän annettuani \VUgras{tarjoilijalle}
\textsuperscript{(57)} kauniin \VUgras{juomarahan} \textsuperscript{(59)}.

%%K t13-10-1 p
10. --- Kun \VUgras{johtaja} \textsuperscript{(60)} näki minun olevan valmis
lähtemään, hän ehdotti minulle, että menisin parvekkeelle ihailemaan
\VUgras{Valtameren} \textsuperscript{(62)} komeaa panoraamaa. Kaukana alkoi
näkyä \VUgras{kalastusveneitä} \textsuperscript{(63)}, \VUgras{purjealuksia}
\textsuperscript{(64)} ja valtava \VUgras{postilaiva} \textsuperscript{(65)},
jonka musta siluetti piirtyy valkoisia \VUgras{pilviä} \textsuperscript{(67)}
vasten \VUgras{taivaanrannalla} \textsuperscript{(66)}. Kevyt tuulahdus sai
liehumaan \VUgras{lipun} \textsuperscript{(70)}, joka on kiinnitetty
\VUgras{kyltin} \textsuperscript{(69)} yläpuolelle, joka on \VUgras{hallissa}
\textsuperscript{(68)}.

%%K t13-tit-3 sub
\VUsaut{3.0mm}
\VUpk{10.2pt}{40}{Huvituksia.}
\VUsaut{3.0mm}

%%K t13-11-1 p
11. --- Näytelmä oli niin kiehtova, että päätin mennä huomenna kahvilan
terassille \VUgras{teatterikiikarin} \textsuperscript{(71)} tai
\VUgras{kaukoputken} \textsuperscript{(72)} kanssa.

%%K t13-12-1 p
12. --- Kun jätin parvekkeen, menin hotellin kahvilaan nielaisemaan
\VUgras{juoman} \textsuperscript{(73)}, samalla kun pelasin
\VUgras{biljardi}-erän \textsuperscript{(75)}. Pysähtymättä kuljin
kahvilahuoneen läpi, jossa monet vieraat pelasivat \VUgras{shakkia}
\textsuperscript{(78)}, \VUgras{tammea} \textsuperscript{(79)},
\VUgras{dominoa} \textsuperscript{(80)}, \VUgras{backgammonia}
\textsuperscript{(81)} tai \VUgras{korttia} \textsuperscript{(82)}, ja menin
suoraan ylös \VUgras{biljardisaliin} \textsuperscript{(74)}.

%%K t13-13-1 p
13. --- Kun erä oli lopussa, laskeuduin pohjakerrokseen ja pyysin tarjoilijalta
\VUgras{kirjekuoren} \textsuperscript{(84)}, \VUgras{paperia}
\textsuperscript{(83)}, \VUgras{postimerkin} \textsuperscript{(85)},
\VUgras{kynän} \textsuperscript{(87)} ja \VUgras{mustetta}
\textsuperscript{(86)} kirjoittaakseni kirjeen.

%%K t13-13-2 p
Sitten kutsuin \VUgras{vuokra-ajurin} \textsuperscript{(92)}, jonka
\VUgras{vaunut} \textsuperscript{(88)} olivat pysähtyneet kahvilan eteen. Kysyin
häneltä hintaa tunnilta tai ajolta, ja, kun olimme hinnasta yksimielisiä, hän
avasi minulle \VUgras{vaununoven} \textsuperscript{(89)} ja asetti minut
paikoilleni. Hän nousi taas \VUgras{ajopukille} \textsuperscript{(91)}, pani
hevoset nopeasti liikkeelle \VUgras{pienellä ruoskalla} \textsuperscript{(93)},
ja lähdimme matkaan hurmaavalle retkelle.
\VUsaut{6.0mm}
\VUfilet{22mm}
